% Extensión estimada: algo más de una página y la tabla de experimentos (ficha 3.9).
% Fuentes: Secciones 3.3 a 3.8 de este capítulo (no agrega experimentos: los ordena por etapas). Ejecuciones por
% configuración: búsqueda con una ejecución por configuración, tres ejecuciones por clasificador supervisado, una por
% función de pérdida; VLM y generadores con decodificación voraz o verosimilitud, deterministas.
% Revisión del tutor del 26/10/2026 (audio 7): la configuración de los experimentos estaba suelta al final de 3.8;
% se ordena por etapas, separando imágenes y reportes, y se dice cuántas ejecuciones tiene cada configuración.
\section{Diseño de los experimentos}
\label{sec:met_diseno}

Las secciones anteriores describen cada componente, el \emph{pipeline} y la forma de evaluarlos. Los experimentos que
los ponen a prueba se organizan por etapas, de modo que cada etapa usa lo que produjo la anterior, y se dividen en dos
partes. La primera parte trabaja con la clasificación de las imágenes y produce el clasificador del
\emph{pipeline}. La segunda trabaja con los reportes y usa ese clasificador como primer componente del
\emph{pipeline}, que solo interviene en esta parte.

La clasificación de las imágenes se evalúa en cuatro etapas:
\begin{enumerate}[leftmargin=1.6em]
    \item \textbf{Búsqueda de hiperparámetros.} Swin-Tiny y ConvNeXt-Tiny se entrenan con las 30 configuraciones de
    la Sección~\ref{sec:met_supervisados_ensamble}, y la validación elige una por modelo.
    \item \textbf{Comparación de los clasificadores supervisados.} Los cuatro modelos se entrenan con la configuración
    elegida, se combinan los dos de la búsqueda en el ensamble y se fija en validación el umbral de melanoma.
    \item \textbf{Variantes y evaluaciones complementarias.} Sobre los modelos de la etapa anterior se comparan las
    funciones de pérdida para el desbalance, se estima el efecto de la fuga de información y se aplican las
    evaluaciones de la Sección~\ref{sec:met_complementarios}.
    \item \textbf{Clasificación con \glspl{vlm}.} Mini-VLM, LLaVA-Med y DermLIP clasifican la misma prueba, sin ajuste y
    con las estrategias y los controles de la Sección~\ref{sec:met_clasif_vlm}, y se comparan con los supervisados.
\end{enumerate}

La generación de reportes, por su parte, avanza en tres etapas, que agregan de a uno los componentes del
\emph{pipeline}:
\begin{enumerate}[leftmargin=1.6em]
    \item \textbf{Generador solo.} Mini-VLM y LLaVA-Med se ajustan sobre \emph{SkinCaRe}
    (Sección~\ref{sec:met_generacion}) y, con SkinVL, generan reportes sin diagnóstico, que se comparan con los textos
    de referencia (Nivel~1) y se evalúan sobre \emph{HAM10000} (Nivel~2).
    \item \textbf{Cascada.} Los tres generadores reciben el diagnóstico de las cuatro fuentes de la
    Sección~\ref{sec:met_referencias}, con las dos variantes que separan la imagen del diagnóstico y con las plantillas
    de recuperación.
    \item \textbf{\emph{Pipeline} completo.} El ruteo por concordancia decide el tipo de reporte de cada caso, y el
    resultado se compara con la cascada y con la línea base sin generador.
\end{enumerate}

El número de ejecuciones depende de si el resultado varía al repetirlo. Cada configuración de la búsqueda y cada
función de pérdida se entrenan una vez, mientras que los cuatro clasificadores supervisados se entrenan tres veces con
la configuración elegida, de modo que sus resultados se informan como media y desvío de las tres ejecuciones
(Sección~\ref{sec:met_seleccion}). Los \glspl{vlm} sin ajuste, las cabezas sobre los \emph{embeddings} y los
generadores, en cambio, se ejecutan una vez por configuración, porque la verosimilitud y la decodificación voraz
devuelven la misma salida para la misma entrada. Cada generador, por su parte, se ajusta una sola vez.

La Tabla~\ref{tab:met_experimentos} reúne los experimentos con lo que busca cada uno, los datos sobre los que se mide y
su métrica principal. El Capítulo~\ref{cp:res_clasif} presenta los resultados de la primera parte, la clasificación
de las imágenes, y el Capítulo~\ref{cp:res_gen}, los de la segunda, la generación de reportes con el
\emph{pipeline}.

\begingroup
\small
\setlength{\tabcolsep}{3pt}
\begin{longtable}{@{}>{\raggedright\arraybackslash}p{3.5cm}
    >{\raggedright\arraybackslash\hspace{0pt}}p{4.3cm}
    >{\raggedright\arraybackslash\hspace{0pt}}p{2.6cm}
    >{\raggedright\arraybackslash\hspace{0pt}}p{3.5cm}@{}}
    \caption{Experimentos del trabajo: qué busca cada uno, sobre qué datos se mide y con qué métrica.}
    \label{tab:met_experimentos}\\
    \toprule
    \textbf{Experimento} & \textbf{Qué se busca} & \textbf{Datos} & \textbf{Métrica principal} \\
    \midrule
    \endfirsthead
    \toprule
    \textbf{Experimento} & \textbf{Qué se busca} & \textbf{Datos} & \textbf{Métrica principal} \\
    \midrule
    \endhead
    \bottomrule
    \endlastfoot
    Clasificadores supervisados (Sección~\ref{sec:met_supervisados})
        & Comparar las dos familias de arquitecturas con el mismo procedimiento
        & Prueba de \emph{HAM10000}
        & Sensibilidad macro; F1 macro \\
    \filasep
    Funciones de pérdida para el desbalance
        & Ver si otra función de pérdida mejora las clases minoritarias
        & Prueba de \emph{HAM10000}
        & Sensibilidad macro \\
    \filasep
    Detección de melanoma
        & Detectar melanomas con un tope de falsas alarmas
        & Prueba de \emph{HAM10000}
        & Sensibilidad y especificidad con umbral; AUC \\
    \filasep
    Fuga de información (Sección~\ref{sec:met_particiones})
        & Estimar cuánto sobreestimaba la fuga el desempeño
        & Prueba completa y subconjunto sin lesiones compartidas
        & Sensibilidad macro \\
    \filasep
    Clasificación con \glspl{vlm} (Sección~\ref{sec:met_clasif_vlm})
        & Saber si un \gls{vlm} clasifica tan bien como un supervisado
        & Prueba de \emph{HAM10000}; fotografías clínicas de \emph{SkinCAP} como control
        & Sensibilidad macro; McNemar con Holm \\
    \filasep
    Generación con referencia
        & Medir el parecido del reporte con un texto de referencia
        & Prueba de \emph{MM-Skin} y de \emph{SkinCoT}
        & Nivel~1 (BLEU-4, ROUGE-L, METEOR, CIDEr, BERTScore) \\
    \filasep
    Ajuste del modelo de lenguaje de Mini-VLM (Sección~\ref{sec:met_generacion})
        & Ver si ajustar también el modelo de lenguaje cambia el reporte
        & Prueba de \emph{HAM10000} y de \emph{SkinCoT}
        & Diferencia pareada de la consistencia diagnóstica; Nivel~1 \\
    \filasep
    Cascada con cuatro fuentes (Sección~\ref{sec:met_referencias})
        & Saber si el reporte afirma el diagnóstico correcto y cuánto aporta el anclaje
        & Prueba de \emph{HAM10000}
        & Consistencia diagnóstica \\
    \filasep
    Cascada sin imagen y con diagnóstico equivocado (Sección~\ref{sec:met_referencias})
        & Separar lo que el reporte toma de la imagen de lo que toma del diagnóstico
        & 400 imágenes de la prueba de \emph{HAM10000}; la prueba completa con SkinVL
        & Consistencia diagnóstica \\
    \filasep
    Recuperación (Sección~\ref{sec:met_recuperacion})
        & Ver si los casos parecidos mejoran el reporte
        & Prueba de \emph{HAM10000}
        & Consistencia diagnóstica \\
    \filasep
    Ruteo en tres niveles (Sección~\ref{sec:met_ruteo})
        & Ver si la concordancia separa aciertos de errores y cuánto aporta el generador
        & Prueba de \emph{HAM10000}
        & Consistencia diagnóstica por nivel; proporción de casos en cada nivel \\
    \filasep
    Generalización (Sección~\ref{sec:met_hiba})
        & Ver si el desempeño se sostiene en otra población
        & \emph{HIBA}
        & Sensibilidad macro \\
    \filasep
    Robustez (Sección~\ref{sec:met_robustez})
        & Medir la caída ante peores condiciones de captura
        & Prueba de \emph{HAM10000} perturbada
        & Caída relativa de la sensibilidad macro \\
    \filasep
    Calibración (Sección~\ref{sec:met_calibracion})
        & Saber cuánto se puede confiar en la probabilidad
        & Validación y prueba de \emph{HAM10000}
        & ECE; Brier \\
    \filasep
    Interpretabilidad (Sección~\ref{sec:met_interpretabilidad})
        & Ver en qué parte de la imagen se apoya el modelo
        & Siete imágenes de prueba
        & Inspección cualitativa \\
    \filasep
    Costo computacional (Sección~\ref{sec:met_costo})
        & Estimar el costo de ejecutar los modelos y el \emph{pipeline}
        & Prueba de \emph{HAM10000}
        & Parámetros; tiempo por imagen y por reporte \\
\end{longtable}
\endgroup
